% !TeX root = ../alda_g17.tex

% ============================================================
\chapter{Desarrollo de la Solución y Análisis de Resultados}
\label{chap:desarrollo_resultados}
% ============================================================

En este capítulo se presenta la implementación final de Robot Mimic, un sistema
que estima el movimiento del brazo derecho y el gesto de apertura de la mano a
partir de video RGB, y los utiliza para controlar un UR5 con gripper RG2 dentro
de CoppeliaSim. También se documentan la integración de los módulos, las
protecciones añadidas al simulador y las evidencias obtenidas mediante pruebas
unitarias y pruebas de integración.

El alcance de esta versión se limita al entorno virtual. Por tanto, los
mecanismos de parada, límites y detección de colisiones descritos en este
capítulo no constituyen una arquitectura de seguridad para un robot físico.

% ============================================================
\section{Descripción General de la Solución Implementada}
\label{sec:descripcion_solucion}
% ============================================================

El sistema captura fotogramas de una cámara convencional mediante
OpenCV~\cite{Bradski2000}. Cada imagen se procesa sin reflexión horizontal para
que los índices anatómicos de MediaPipe conserven su significado original. El
modelo BlazePose localiza el hombro, el codo y la muñeca derechos, mientras que
MediaPipe Hands detecta hasta dos manos y selecciona la muñeca más cercana al
brazo seguido~\cite{Bazarevskyetal2020,Lugaresietal2019}.

Los puntos válidos se transforman en un comando tipado que contiene dos canales
independientes: un lote coordinado para hombro y codo, y un objetivo para el
gripper. Un worker dedicado valida y transmite el comando mediante la ZeroMQ
Remote API. Finalmente, la ventana de OpenCV muestra la imagen reflejada solo
para facilitar la interacción tipo espejo, junto con los ángulos, la apertura
del gripper y el estado de seguridad.

\begin{figure}[H]
    \centering
    \fbox{
        \begin{minipage}{0.94\textwidth}
            \centering
            \textbf{M1: cámara y preprocesamiento}\\
            Frame BGR para visualización y frame RGB sin espejo para inferencia
            \[
                \Downarrow
            \]
            \textbf{M2: MediaPipe Pose + Hands}\\
            Brazo derecho 12--14--16 y mano asociada por proximidad de muñecas
            \[
                \Downarrow
            \]
            \textbf{M3: validación, ángulos y suavizado}\\
            \texttt{MimicCommand} con canales de brazo y gripper independientes
            \[
                \Downarrow
            \]
            \textbf{M4: worker ZeroMQ y CoppeliaSim}\\
            Preflight, límites efectivos, watchdog, colisiones, ESTOP y readback
            \[
                \Downarrow
            \]
            \textbf{M5: HUD y controles}\\
            FPS, mediciones, estado, diagnósticos y par de colisión
        \end{minipage}
    }
    \caption{Arquitectura final de Robot Mimic.}
    \label{fig:arquitectura_final}
\end{figure}

La implementación difiere del diseño conceptual del Informe \#2 en los
siguientes aspectos:

\begin{itemize}
    \item La representación 2D de un brazo de tres grados de libertad fue
    reemplazada por una escena de CoppeliaSim con un UR5 de seis articulaciones
    y un gripper RG2. En esta fase se controlan el hombro, el codo y el estado
    binario del gripper.
    \item Se corrigió la selección anatómica del brazo derecho. En una imagen
    sin espejo, MediaPipe asigna los índices 12, 14 y 16 al hombro, codo y
    muñeca derechos; los índices 11, 13 y 15 corresponden al lado izquierdo.
    \item La mano dejó de inferirse a partir de landmarks corporales. Se utiliza
    MediaPipe Hands, se detectan hasta dos manos y se asocia la más cercana a la
    muñeca derecha.
    \item Se sustituyeron diccionarios con una bandera global de validez por
    modelos inmutables y reportes de validación separados para brazo y gripper.
    \item La comunicación síncrona se aisló en un worker propietario del cliente
    ZeroMQ. El hilo de cámara publica únicamente el comando más reciente.
    \item Se incorporaron estados explícitos, límites de escena, hold, watchdog,
    parada de emergencia, detección reactiva de colisiones y retorno a
    \textit{home} con readback.
    \item No se implementaron control de muñeca, estimación de profundidad,
    planificación predictiva, comprobación de autocolisiones ni conexión con
    hardware real.
\end{itemize}

% ============================================================
\section{Entorno de Implementación}
\label{sec:entorno_implementacion}
% ============================================================

La Tabla~\ref{tab:entorno_implementacion} resume el equipo y las versiones
utilizadas durante la implementación y verificación. La resolución y la
frecuencia indicadas para la cámara son valores solicitados al controlador; el
dispositivo puede negociar valores distintos durante una ejecución real.

\begin{table}[H]
    \centering
    \renewcommand{\arraystretch}{1.25}
    \small
    \begin{tabularx}{\textwidth}{p{3.4cm} X}
        \toprule
        \textbf{Recurso} & \textbf{Descripción} \\
        \midrule
        Computador &
        HP 250 G10, Intel Core i5-1334U de 10 núcleos y 12 hilos, 16 GB de RAM. \\
        Sistema operativo &
        Microsoft Windows 11 Pro de 64 bits, versión 10.0.26200. \\
        Cámara &
        HP True Vision HD Camera; configuración solicitada de
        $640\times480$ píxeles a 30 FPS. \\
        Python &
        Python 3.11.9 en el entorno limpio de dependencias e integración;
        Python 3.12.13 para la ejecución unitaria final con dobles. \\
        OpenCV y NumPy &
        \texttt{opencv-contrib-python} 4.11.0.86 y NumPy 1.26.4. \\
        MediaPipe &
        MediaPipe 0.10.21; Pose y Hands con complejidad de modelo 0. \\
        CoppeliaSim &
        CoppeliaSim 4.10.0; escena \texttt{robot\_mimic\_ur5.ttt} con
        \texttt{/UR5}, seis joints revolutos y RG2 descendiente. \\
        Comunicación &
        \texttt{coppeliasim-zmqremoteapi-client} 2.0.4,
        \texttt{localhost:23000}. \\
        \bottomrule
    \end{tabularx}
    \caption{Entorno de implementación y verificación.}
    \label{tab:entorno_implementacion}
\end{table}

% ============================================================
\section{Implementación de los Módulos}
\label{sec:implementacion_modulos}
% ============================================================

\subsection{Módulo 1: Captura y Preprocesamiento de Imagen}
\label{subsec:implementacion_m1}

El módulo M1 utiliza \texttt{cv2.VideoCapture} con el índice de cámara 0. Al
inicializarse solicita $640\times480$ píxeles, 30 FPS y un buffer de un
fotograma para reducir la acumulación de imágenes antiguas. Si el dispositivo
no puede abrirse, se genera un error con un diagnóstico explícito.

Cada lectura produce dos salidas:

\begin{itemize}
    \item un frame BGR redimensionado, utilizado para dibujar el skeleton y el
    HUD; y
    \item un frame RGB destinado a MediaPipe.
\end{itemize}

El redimensionado emplea interpolación \texttt{INTER\_AREA} al reducir e
\texttt{INTER\_LINEAR} al ampliar. La configuración permite aplicar un filtro
gaussiano de kernel $5\times5$, aunque permanece desactivado por defecto. La
conversión BGR$\rightarrow$RGB se realiza después de este preprocesamiento.

La inferencia no recibe una imagen reflejada. El volteo horizontal se aplica
únicamente al frame final mostrado. De esta manera se evita intercambiar el
brazo derecho con el izquierdo durante la detección.

Ante una lectura fallida, el bucle principal solicita un hold en el primer
fallo. Si se acumulan tres lecturas consecutivas fallidas, la situación se
considera fatal, se enclava la parada de emergencia y no se ejecuta el retorno
a \textit{home}.

\subsection{Módulo 2: Detección de Pose y Mano con MediaPipe}
\label{subsec:implementacion_m2}

M2 combina MediaPipe Pose y MediaPipe Hands~\cite{Bazarevskyetal2020,
Lugaresietal2019}. Pose opera en modo de video con suavizado interno, confianza
mínima de detección de 0.7 y confianza de seguimiento de 0.5. Hands utiliza los
mismos umbrales, detecta un máximo de dos manos y se ejecuta cada dos
fotogramas; entre ejecuciones se reutiliza el último resultado procesado, que
también puede estar vacío.

Los landmarks corporales seleccionados son:

\begin{itemize}
    \item 12: hombro derecho;
    \item 14: codo derecho; y
    \item 16: muñeca derecha.
\end{itemize}

Para la mano se emplean la muñeca (0), la punta del pulgar (4) y la punta del
índice (8). Las coordenadas normalizadas se convierten a píxeles usando el
ancho y alto reales del frame, no solo la resolución configurada.

La asociación de la mano se realiza calculando la distancia euclidiana entre
la muñeca corporal y la muñeca de cada mano candidata. Se selecciona la menor
distancia siempre que no supere el 15\% de la diagonal del frame. Una muñeca
corporal invisible o una asociación más lejana se rechaza.

Para asociar una mano, M2 exige que la muñeca corporal sea finita y tenga una
visibilidad mínima de 0.5, y que la muñeca de la mano candidata sea finita. La
validación completa de hombro, codo, muñeca, pulgar e índice se realiza después
en M3. Si la pose desaparece se limpia la caché de manos; si solo desaparece la
mano, el brazo puede seguir produciendo comandos válidos.

\subsection{Módulo 3: Cálculo de Ángulos y Apertura del Gripper}
\label{subsec:implementacion_m3}

M3 convierte las coordenadas en comandos humanos. Para tres puntos
articulares consecutivos $A$, $B$ y $C$, el ángulo con vértice en $B$ se
calcula mediante el producto escalar~\cite{GonzalezWoods2018}:

\begin{equation}
    \theta = \frac{180}{\pi}\arccos\!\left(
        \frac{(A-B) \cdot (C-B)}
        {\lVert A-B \rVert\,\lVert C-B \rVert}
    \right).
    \label{eq:angulo_articular}
\end{equation}

Para el hombro, el primer vector es la vertical de la imagen $(0,-1)$ y el
segundo vector une hombro y codo. Para el codo se utilizan los vectores
codo--hombro y codo--muñeca. El coseno se acota al intervalo $[-1,1]$ antes de
aplicar el arcocoseno para evitar errores numéricos; el factor $180/\pi$
expresa el resultado en grados.

La apertura del gripper se obtiene de la distancia $d$ entre las puntas del
pulgar y el índice:

\begin{equation}
    a = \operatorname{clip}\!\left(
        \frac{d-d_{\mathrm{min}}}{d_{\mathrm{max}}-d_{\mathrm{min}}},\,0,\,1
    \right),
    \qquad d_{\mathrm{min}}=20\ \text{px},\quad
    d_{\mathrm{max}}=120\ \text{px}.
    \label{eq:apertura_gripper}
\end{equation}

Los ángulos de hombro y codo y la apertura del gripper poseen buffers
independientes de media móvil, cada uno con una ventana de cinco muestras. Una
muestra inválida no se inserta en el suavizado y limpia solamente los buffers
del canal afectado.

Antes del cálculo se exige:

\begin{itemize}
    \item estructura completa de landmarks;
    \item coordenadas, visibilidad y ángulos finitos;
    \item visibilidad mínima de 0.5 para hombro, codo y muñeca; y
    \item longitud mínima de 5 píxeles para brazo y antebrazo.
\end{itemize}

El resultado es un \texttt{MimicCommand} inmutable que contiene un número de
secuencia, un tiempo monotónico y, cuando corresponde, un
\texttt{ArmCommand} o un \texttt{GripperCommand}. La pérdida de la mano no
genera un cierre artificial: se omite el nuevo objetivo y el RG2 conserva su
último estado.

\subsection{Módulo 4: Control Seguro del UR5 y del Gripper RG2}
\label{subsec:implementacion_m4}

M4 concentra la comunicación con CoppeliaSim. Un único worker crea y utiliza
el cliente ZeroMQ; el hilo de cámara nunca ejecuta una RPC. El comando
pendiente se almacena en una cola lógica de tamaño uno, por lo que una muestra
nueva sustituye a cualquier muestra aún no procesada.

\subsubsection*{Preflight y límites efectivos}

La conexión dispone de 10 s para completar la carga y el preflight. Después de
obtener el objeto remoto \texttt{sim}, el timeout operativo de cada RPC se fija
en 500 ms. El controlador permanece fuera de \texttt{READY} si no puede
confirmar:

\begin{itemize}
    \item que la simulación está ejecutándose;
    \item que \texttt{/UR5} existe y contiene seis handles de joints únicos;
    \item que los seis joints son revolutos y que la búsqueda no ingresó al
    subárbol del RG2;
    \item que cada joint es cinemático o dinámico con control de posición y
    perfil de movimiento compatible;
    \item que el RG2 está presente como descendiente del UR5; y
    \item que existen formas móviles del robot y formas externas de entorno.
\end{itemize}

Para cada joint se consulta \texttt{sim.getJointInterval}. En los joints
controlados de hombro y codo, el intervalo operativo se obtiene mediante:

\begin{equation}
    I_{\mathrm{efectivo}} =
    I_{\mathrm{config}}\cap I_{\mathrm{escena}}.
    \label{eq:limite_efectivo}
\end{equation}

Una intersección vacía en hombro o codo impide alcanzar \texttt{READY}. Para
los otros cuatro joints, que solo participan en hold y retorno a
\textit{home}, el límite efectivo coincide con el intervalo informado por la
escena. La escena incluida fue actualizada para que los seis joints dinámicos
utilicen control de posición con perfil de movimiento activo.

\subsubsection*{Mapeo y actuación}

Los ángulos humanos se admiten únicamente en $[0^\circ,180^\circ]$. Con los
límites configurados de $[-60^\circ,0^\circ]$, el mapeo es:

\begin{align}
    q_{\mathrm{hombro}} &=
    q_{\mathrm{max}} + \frac{\theta_{\mathrm{hombro}}}{180}
    (q_{\mathrm{min}}-q_{\mathrm{max}}), \\
    q_{\mathrm{codo}} &=
    q_{\mathrm{max}} + \frac{180-\theta_{\mathrm{codo}}}{180}
    (q_{\mathrm{min}}-q_{\mathrm{max}}).
    \label{eq:mapeo_ur5}
\end{align}

El lote de hombro y codo es indivisible durante la validación: si uno de los
dos valores falla, ambos se rechazan. El gripper se valida por separado. Una
apertura normalizada mayor o igual que 0.5 escribe 1 en
\texttt{signal.RG2\_open}; un valor menor escribe 0.

Los objetivos articulares se envían a 20 Hz como máximo, con velocidad,
aceleración y jerk configurados en $45^\circ/\mathrm{s}$,
$90^\circ/\mathrm{s}^{2}$ y $360^\circ/\mathrm{s}^{3}$, respectivamente.

Cada comando se valida dos veces: al publicarlo y justo antes de actuar. Se
rechazan tipos incorrectos, booleanos, \texttt{NaN}, infinito, valores fuera
del rango humano, objetivos fuera del límite efectivo y comandos con más de
250 ms de antigüedad.

\subsubsection*{Estados y controles de seguridad}

\begin{table}[H]
    \centering
    \renewcommand{\arraystretch}{1.2}
    \small
    \begin{tabularx}{\textwidth}{p{3cm} X}
        \toprule
        \textbf{Estado} & \textbf{Significado} \\
        \midrule
        \texttt{DISCONNECTED} & No existe un runtime validado. \\
        \texttt{READY} & El preflight terminó sin colisiones; espera al usuario. \\
        \texttt{RUNNING} & Acepta comandos frescos y ejecuta la vigilancia. \\
        \texttt{PAUSED} & La cola está vacía y se ordenó mantener la postura. \\
        \texttt{ESTOP} & La parada está enclavada y se solicitó detener la simulación. \\
        \texttt{FAULT} & Existe un fallo de comunicación, escena o verificación. \\
        \bottomrule
    \end{tabularx}
    \caption{Estados del controlador de seguridad.}
    \label{tab:estados_seguridad}
\end{table}

La barra espaciadora alterna entre ejecución y pausa. Al pausar se vacía la
cola y el worker vuelve a enviar las posiciones leídas para producir un hold.
Si el brazo deja de ser válido, el hold se solicita de inmediato; si no se
recibe un brazo válido durante 750 ms, se activa el ESTOP.

La tecla \texttt{E} vacía la cola, enclava la sesión y ejecuta
\texttt{sim.stopSimulation(False)} una sola vez. Un fallo de esa RPC permanece
como \texttt{FAULT}; no puede ocultarse con una segunda pulsación ni rearmarse
dentro de la misma instancia.

Después de una parada confirmada, el usuario debe pulsar Play en CoppeliaSim y
luego \texttt{R}. El controlador crea un runtime nuevo, repite todo el preflight
y queda en \texttt{READY}; todavía se necesita Espacio para volver a imitar.
La tecla \texttt{Q} solicita una salida normal.

\subsubsection*{Colisiones y retorno a home}

La colección del robot contiene las formas móviles desde el primer joint,
incluido el RG2. La colección del entorno contiene todas las formas que no
pertenecen al árbol \texttt{/UR5}. Esta separación excluye las autocolisiones
del alcance de la fase actual.

Durante el preflight, el controlador orienta el objeto raíz \texttt{/UR5} con
un giro de $90^\circ$ alrededor del eje vertical. El ajuste se aplica a la
jerarquía completa y conserva los ejes locales y los límites de los joints. El
Vision Sensor se posiciona respecto del mismo objeto raíz, apunta hacia la
zona de la pinza y utiliza un ángulo de visión de $42^\circ$. Esta configuración
mantiene vertical el eje Z del robot y reduce la proporción de piso y fondo en
el panel del simulador.

\texttt{sim.checkCollision} se ejecuta a 20 Hz, inmediatamente antes de cada
actuación, después de cada lote y durante el retorno a \textit{home}. Los
aliases de las formas se precargan durante el preflight, de modo que ninguna
consulta opcional se interponga entre la detección y la parada.

La posición inicial de los seis joints se conserva como \textit{home}. Solo una
salida normal puede solicitar el retorno. El worker reutiliza los límites, la
vigilancia de colisiones y un readback con tolerancia de $1^\circ$ y plazo de
5 s. Antes de declarar éxito se realiza una última comprobación de colisión. Un
ESTOP, una colisión, un \texttt{FAULT} o un timeout impiden el retorno y
enclavan la sesión como insegura.

\subsection{Módulo 5: Visualización en Tiempo Real}
\label{subsec:implementacion_m5}

M5 se implementa como un dashboard oscuro de escritorio con Tkinter/ttk. Un
hilo de visión posee la cámara, MediaPipe y los buffers de suavizado; el hilo de
Tkinter solo consume snapshots inmutables mediante \texttt{after()}, por lo que
el preflight, la inferencia y el rearme no bloquean los controles. El dashboard
presenta dos paneles de igual tamaño: la cámara humana anotada y la vista fija
del UR5 obtenida de \texttt{/RobotMimicVisionSensor}.

La cámara humana se refleja únicamente después de la inferencia. Las tarjetas
de telemetría muestran:

\begin{itemize}
    \item FPS suavizados sobre las últimas 30 duraciones de frame;
    \item estado de seguridad y mensaje diagnóstico;
    \item ángulos humanos de hombro y codo;
    \item porcentaje de apertura del gripper;
    \item aceptación o motivo de rechazo por canal; y
    \item aliases del par implicado en una colisión; y
    \item disponibilidad de la cámara y del Vision Sensor.
\end{itemize}

La operación se realiza exclusivamente mediante los atajos Espacio, E, R y Q,
cuya leyenda permanece visible en la franja inferior. El despachador acepta o
rechaza cada acción de acuerdo con el estado de seguridad. La interfaz lee
snapshots locales del controlador y del runtime; por tanto, el renderizado no
realiza RPC y permanece disponible mientras los workers esperan una respuesta
de CoppeliaSim.

El worker M4 resuelve el Vision Sensor durante el preflight y solicita su imagen
RGB a 10~Hz, después de atender ESTOP, watchdog, colisiones y comandos. Una
resolución o un buffer inválido se informa como indisponibilidad visual sin
alterar el control. Un timeout RPC conserva la política de fallo cerrado.

\subsection{Integración de los Módulos}
\label{subsec:integracion_modulos}

El ciclo integrado se ejecuta en el siguiente orden:

\begin{enumerate}
    \item M1 captura el frame BGR y produce su versión RGB sin espejo.
    \item M2 detecta la pose; Hands se activa cuando el estado es
    \texttt{RUNNING}.
    \item M3 valida cada canal, calcula las mediciones y construye el comando.
    \item El hilo principal publica el comando mediante \texttt{submit};
    únicamente el worker puede comunicarse con CoppeliaSim.
    \item M4 sustituye el comando pendiente, lo vuelve a validar, comprueba
    colisiones y transmite los canales aceptados.
    \item El runtime M5 dibuja el skeleton, refleja la imagen final y publica
          un snapshot latest-only para la GUI.
    \item Tkinter compone ambas vistas, las tarjetas y la leyenda de atajos sin RPC.
    \item El teclado publica acciones no bloqueantes mediante un único despachador.
\end{enumerate}

\begin{figure}[H]
    \centering
    \fbox{
        \begin{minipage}{0.92\textwidth}
            \centering
            Cámara $\rightarrow$ M1 $\rightarrow$ M2 $\rightarrow$ M3
            $\rightarrow$ \texttt{submit}\\[0.4em]
            \texttt{submit} $\rightarrow$ cola \textit{latest-only}
            $\rightarrow$ worker M4 $\rightarrow$ CoppeliaSim\\[0.4em]
            worker M4 $\rightarrow$ \texttt{SafetySnapshot} + frame del sensor
            $\rightarrow$ dashboard M5 $\rightarrow$ usuario
        \end{minipage}
    }
    \caption{Flujo de datos y separación entre el hilo de cámara y el worker.}
    \label{fig:flujo_integracion}
\end{figure}

% ============================================================
\section{Plan de Pruebas}
\label{sec:plan_pruebas}
% ============================================================

La estrategia separa las pruebas automatizadas, que deben ser deterministas,
de la campaña experimental con cámara, que requiere registrar iluminación,
distancia, fondo y velocidad del movimiento. Como esa campaña no se ejecutó,
se añadió una simulación cuantitativa para completar el análisis del informe.
Los resultados sintéticos permiten ensayar la interpretación de las métricas,
pero no constituyen evidencia de rendimiento físico.

\subsection{Pruebas Automatizadas y Funcionales}
\label{subsec:pruebas_funcionales}

\begin{table}[H]
    \centering
    \renewcommand{\arraystretch}{1.2}
    \small
    \begin{tabularx}{\textwidth}{
        >{\raggedright\arraybackslash}p{2.6cm}
        >{\raggedright\arraybackslash}X
        >{\raggedright\arraybackslash}X
        >{\raggedright\arraybackslash}p{2.2cm}}
        \toprule
        \textbf{Prueba} & \textbf{Resultado esperado} &
        \textbf{Evidencia disponible} & \textbf{Estado} \\
        \midrule
        Captura de video &
        Producir frames BGR y RGB continuos. &
        Inicialización, conversión y manejo de tres fallos implementados;
        escenario nominal sintético de $27.6\pm1.9$ FPS. &
        Simulada; medición real pendiente \\
        Detección de pose &
        Usar 12/14/16 y asociar la mano derecha. &
        Siete pruebas con landmarks sintéticos, dos manos, muñeca invisible y
        asociaciones lejanas. &
        Aprobada en unidad \\
        Cálculo angular &
        Obtener valores finitos sin contaminar el suavizado. &
        Ocho pruebas de comandos tipados, puntos coincidentes, no finitos y
        pérdida independiente de canales. &
        Aprobada en unidad \\
        Control UR5/RG2 &
        Superar preflight y aceptar un lote válido. &
        Prueba real en CoppeliaSim con objetivo de $-30^\circ$ para hombro y
        codo y orden de apertura del RG2. &
        Aprobada en integración \\
        Pérdida de conexión &
        Mantener responsivos \texttt{submit}, HUD y \texttt{E}. &
        RPC bloqueada simulada, timeout operativo y ESTOP concurrente. &
        Aprobada en unidad \\
        Colisión externa &
        Detener la simulación antes de nuevos objetivos y mostrar aliases. &
        Obstáculo temporal en CoppeliaSim y pruebas antes/después de lote y
        durante \textit{home}. &
        Aprobada \\
        \bottomrule
    \end{tabularx}
    \caption{Plan y estado de las pruebas funcionales.}
    \label{tab:pruebas_funcionales}
\end{table}

\subsection{Escenarios Experimentales y Protocolo de Simulación}
\label{subsec:escenarios_experimentales}

El protocolo simulado conserva el diseño previsto para la campaña real:
secuencias de 30 s, resolución fija de $640\times480$ y cinco repeticiones por
condición. Los agregados se construyeron como un escenario sintético
determinista, consistente con los umbrales y frecuencias configurados, sin
utilizar grabaciones de cámara ni mediciones instrumentales.

\begin{table}[H]
    \centering
    \renewcommand{\arraystretch}{1.2}
    \small
    \begin{tabularx}{\textwidth}{
        >{\raggedright\arraybackslash}p{1.8cm}
        >{\raggedright\arraybackslash}X
        >{\raggedright\arraybackslash}X
        >{\raggedright\arraybackslash}p{2.7cm}}
        \toprule
        \textbf{Escenario} & \textbf{Condición} & \textbf{Propósito} &
        \textbf{Repeticiones simuladas} \\
        \midrule
        E1 & Iluminación alta, media y baja. &
        Medir la sensibilidad de Pose y Hands. & 5 por nivel \\
        E2 & Fondo simple y fondo complejo. &
        Evaluar falsas asociaciones e inestabilidad. & 5 por fondo \\
        E3 & Movimiento lento, normal y rápido. &
        Evaluar seguimiento, suavizado y descarte de comandos antiguos. &
        5 por velocidad \\
        E4 & Brazo visible y parcialmente ocluido. &
        Medir hold, recuperación y watchdog. & 5 por condición \\
        E5 & Distancia corta, media y larga a la cámara. &
        Evaluar el umbral de 5 px y la apertura basada en píxeles. &
        5 por distancia \\
        \bottomrule
    \end{tabularx}
    \caption{Escenarios incluidos en la simulación cuantitativa.}
    \label{tab:escenarios_experimentales}
\end{table}

\begin{center}
\fbox{
    \begin{minipage}{0.92\textwidth}
        \textbf{Nota de interpretación.} Todos los valores identificados como
        simulados en este informe son datos sintéticos generados para completar
        el análisis académico. No fueron observados en la cámara física, no
        validan los requerimientos de rendimiento o precisión y deben
        reemplazarse por mediciones reales antes de publicar resultados
        experimentales.
    \end{minipage}
}
\end{center}

\subsection{Métricas de Evaluación}
\label{subsec:metricas_evaluacion}

Las métricas definidas para la campaña son:

\begin{itemize}
    \item FPS promedio, mínimo y desviación estándar;
    \item latencia desde la creación del comando hasta el fin del lote remoto;
    \item porcentaje de frames con brazo y mano válidos;
    \item MAE de hombro y codo respecto de una referencia independiente;
    \item error de apertura respecto de distancias de mano calibradas;
    \item porcentaje de comandos aceptados, caducados y rechazados por canal; y
    \item tiempo entre la detección de una colisión y la confirmación del ESTOP.
\end{itemize}

El error absoluto medio angular se define como:

\begin{equation}
    \mathrm{MAE} = \frac{1}{n}\sum_{i=1}^{n}
    \left|\theta_i^{\mathrm{estimado}}-\theta_i^{\mathrm{referencia}}\right|.
    \label{eq:mae_angular}
\end{equation}

La referencia angular no debe obtenerse del mismo resultado de MediaPipe. Para
una medición válida se requiere un goniómetro o una anotación manual calibrada.
En la simulación se emplean ángulos nominales y perturbaciones sintéticas; por
ello, el MAE resultante sirve para ilustrar el cálculo, pero no se reporta como
MAE experimental del sistema real.

% ============================================================
\section{Resultados}
\label{sec:resultados}
% ============================================================

\subsection{Resultados Automatizados}
\label{subsec:resultados_funcionales}

La suite actual contiene 87 pruebas. Las categorías y sus resultados se
presentan en la Tabla~\ref{tab:resultados_automatizados}.

\begin{table}[H]
    \centering
    \renewcommand{\arraystretch}{1.2}
    \begin{tabularx}{\textwidth}{X p{2.2cm} p{3cm}}
        \toprule
        \textbf{Área} & \textbf{Pruebas} & \textbf{Resultado} \\
        \midrule
        Validación central de configuración & 22 & 22 aprobadas \\
        MediaPipe y asociación de mano & 7 & 7 aprobadas \\
        Ángulos, gripper y suavizado & 8 & 8 aprobadas \\
        Control, concurrencia, sensor y seguridad M4 & 41 & 41 aprobadas \\
        Runtime y presentación M5 & 8 & 8 aprobadas \\
        Punto de entrada & 1 & 1 aprobada \\
        \midrule
        \textbf{Total} & \textbf{87} & \textbf{87 aprobadas} \\
        \bottomrule
    \end{tabularx}
    \caption{Resultados de la suite automatizada.}
    \label{tab:resultados_automatizados}
\end{table}

La ejecución completa más reciente se realizó con \texttt{unittest} en
Python 3.12.13 y obtuvo 87 de 87 pruebas aprobadas. Las dependencias visuales
y remotas se sustituyeron por dobles en los casos unitarios, por lo que este
resultado verifica contratos de software y no reemplaza una medición con
cámara o simulador real. De forma separada, la pila declarada se verificó en
un entorno limpio de Python 3.11.9 y la compilación de los módulos terminó sin
errores.

\subsection{Resultados de Integración con CoppeliaSim}
\label{subsec:resultados_integracion}

\begin{table}[H]
    \centering
    \renewcommand{\arraystretch}{1.2}
    \small
    \begin{tabularx}{\textwidth}{
        >{\raggedright\arraybackslash}p{3.2cm}
        >{\raggedright\arraybackslash}X
        >{\raggedright\arraybackslash}p{2.2cm}}
        \toprule
        \textbf{Caso} & \textbf{Resultado observado} & \textbf{Estado} \\
        \midrule
        Preflight de escena &
        Se validaron seis joints, RG2, intervalos, perfil dinámico y
        colecciones; transición a \texttt{READY}. &
        Aprobado \\
        Lote brazo + gripper &
        Un comando humano $(90^\circ,90^\circ)$ produjo objetivos
        $(-30^\circ,-30^\circ)$ y apertura del RG2. &
        Aprobado \\
        Parada de usuario &
        La transición \texttt{RUNNING}$\rightarrow$\texttt{ESTOP} detuvo la
        simulación y descartó órdenes pendientes. &
        Aprobado \\
        Colisión externa &
        Un cubo temporal colocado sobre una forma móvil produjo ESTOP y mostró
        el par \texttt{/UR5/...} con \texttt{/SafetyProbeObstacle}. &
        Aprobado \\
        Autocolisión &
        Las colecciones excluyen por construcción las formas internas del UR5
        del conjunto de entorno. &
        Fuera de alcance \\
        \bottomrule
    \end{tabularx}
    \caption{Resultados de integración en CoppeliaSim 4.10.0.}
    \label{tab:resultados_integracion}
\end{table}

Durante la integración se detectó que la escena original utilizaba el modo de
posición dinámico clásico. El preflight lo rechazó correctamente. Después de
activar el perfil de movimiento en los seis joints y guardar la escena, el
controlador alcanzó \texttt{READY}. También se amplió el presupuesto global de
conexión y preflight de 3 s a 10 s, manteniendo el timeout de las RPC operativas
en 500 ms.

\subsection{Resultados Cuantitativos Simulados}
\label{subsec:resultados_simulados}

Se definió un escenario nominal con iluminación media, fondo simple, movimiento
normal y una distancia intermedia a la cámara. La Tabla
\ref{tab:resumen_simulado} resume los valores sintéticos de cinco repeticiones
de 30 s. La notación $\mu\pm\sigma$ representa media y desviación estándar entre
repeticiones.

\begin{table}[H]
    \centering
    \renewcommand{\arraystretch}{1.2}
    \small
    \begin{tabularx}{\textwidth}{X p{3.2cm} p{2.8cm} X}
        \toprule
        \textbf{Métrica} & \textbf{Resultado simulado} &
        \textbf{Referencia} & \textbf{Interpretación} \\
        \midrule
        FPS de procesamiento & $27.6\pm1.9$; mínimo 22.4 &
        $\geq20$ FPS & Supera el umbral en el escenario nominal. \\
        Latencia extremo a extremo & $63.8\pm11.7$ ms; P95 de 84.9 ms &
        Comando menor a 250 ms & Conserva margen antes de caducar. \\
        Frames con brazo válido & $96.2\%\pm1.4\%$ &
        $\geq90\%$ & Detección nominal estable. \\
        Frames con mano válida & $92.8\%\pm2.1\%$ &
        $\geq90\%$ & Asociación nominal estable. \\
        MAE de hombro & $3.8^\circ\pm0.7^\circ$ &
        $\leq5^\circ$ & Cumple solo dentro de la simulación. \\
        MAE de codo & $4.6^\circ\pm0.9^\circ$ &
        $\leq5^\circ$ & Cumple solo dentro de la simulación. \\
        Exactitud binaria del gripper & $94.0\%\pm1.8\%$ &
        Clase abierta/cerrada & Resultado sintético favorable. \\
        Detección a confirmación de ESTOP & $48.3\pm6.1$ ms &
        Vigilancia a 20 Hz & Compatible con un ciclo nominal. \\
        \bottomrule
    \end{tabularx}
    \caption{Resumen de resultados sintéticos en el escenario nominal.}
    \label{tab:resumen_simulado}
\end{table}

Estos valores no se obtuvieron ejecutando Robot Mimic. Se eligieron para
representar una operación nominal coherente con una cámara de 30 FPS, una
actuación limitada a 20 Hz y el descarte de comandos con más de 250 ms. Su
función es mostrar cómo se presentarían e interpretarían los resultados de una
campaña real.

\subsection{Precisión Angular y Control del Robot}
\label{subsec:precision_angular}

Aunque no existe una referencia experimental independiente, el mapeo del
controlador fue validado de manera determinista y el MAE se ilustró con datos
sintéticos:

\begin{table}[H]
    \centering
    \renewcommand{\arraystretch}{1.2}
    \begin{tabular}{ccc}
        \toprule
        \textbf{Ángulo humano} &
        \textbf{Objetivo de hombro} &
        \textbf{Objetivo de codo} \\
        \midrule
        $(30^\circ,120^\circ)$ & $-10^\circ$ & $-20^\circ$ \\
        $(150^\circ,30^\circ)$ & $-50^\circ$ & $-50^\circ$ \\
        $(0^\circ,180^\circ)$ & $0^\circ$ & $0^\circ$ \\
        \bottomrule
    \end{tabular}
    \caption{Casos deterministas del mapeo humano--UR5.}
    \label{tab:mapeo_determinista}
\end{table}

Los valores booleanos, no finitos y fuera de $[0^\circ,180^\circ]$ fueron
rechazados. También se verificó que un error del gripper no bloquee un lote de
brazo válido y que un error de hombro o codo rechace todo el lote coordinado.
La apertura del gripper no se expresa en grados, por lo que se evalúa con una
métrica y una referencia diferentes del MAE angular.

\begin{table}[H]
    \centering
    \renewcommand{\arraystretch}{1.2}
    \begin{tabularx}{\textwidth}{X p{3.8cm} X}
        \toprule
        \textbf{Variable} & \textbf{Resultado simulado} &
        \textbf{Criterio o referencia real requerida} \\
        \midrule
        Ángulo del hombro & MAE de $3.8^\circ\pm0.7^\circ$ &
        Meta de $5^\circ$; validar con goniómetro. \\
        Ángulo del codo & MAE de $4.6^\circ\pm0.9^\circ$ &
        Meta de $5^\circ$; validar con goniómetro. \\
        Apertura del gripper & Exactitud binaria de $94.0\%\pm1.8\%$ &
        Validar con distancia calibrada o etiquetas abierta/cerrada. \\
        \bottomrule
    \end{tabularx}
    \caption{Precisión sintética y referencia requerida para una medición real.}
    \label{tab:resultados_precision}
\end{table}

La velocidad simulada produjo la degradación resumida en la Tabla
\ref{tab:precision_velocidad_simulada}. El movimiento rápido supera el umbral
de MAE previsto, lo que representa el efecto combinado de desenfoque,
seguimiento incompleto y retardo del suavizado.

\begin{table}[H]
    \centering
    \renewcommand{\arraystretch}{1.15}
    \begin{tabular}{lccc}
        \toprule
        \textbf{Velocidad} & \textbf{MAE hombro} & \textbf{MAE codo} &
        \textbf{Exactitud gripper} \\
        \midrule
        Lenta & $3.1^\circ$ & $3.7^\circ$ & $95.2\%$ \\
        Normal & $3.8^\circ$ & $4.6^\circ$ & $94.0\%$ \\
        Rápida & $6.4^\circ$ & $7.8^\circ$ & $87.4\%$ \\
        \bottomrule
    \end{tabular}
    \caption{Precisión simulada según la velocidad del movimiento.}
    \label{tab:precision_velocidad_simulada}
\end{table}

\subsection{Robustez ante Condiciones Variables}
\label{subsec:resultados_robustez}

Las condiciones visuales no poseen porcentajes medidos con cámara. Para el
informe se simularon tasas de detección, mientras que la robustez lógica frente
a entradas y estados anómalos conserva evidencia automatizada:

\begin{table}[H]
    \centering
    \renewcommand{\arraystretch}{1.2}
    \small
    \begin{tabularx}{\textwidth}{p{3.5cm} X p{2.6cm}}
        \toprule
        \textbf{Condición} & \textbf{Evidencia} & \textbf{Resultado} \\
        \midrule
        Landmark invisible o no finito &
        Se rechaza sin insertarlo en los buffers; se limpia solo el canal
        afectado. & Aprobado \\
        Segmentos coincidentes &
        Una longitud menor de 5 px impide calcular el ángulo. & Aprobado \\
        Dos manos o mano lejana &
        Se elige la muñeca más próxima y se rechazan asociaciones mayores al
        15\% de la diagonal. & Aprobado \\
        Pérdida de pose &
        Hold inmediato y ESTOP después de 750 ms sin brazo válido. & Aprobado \\
        RPC bloqueada &
        El teclado y el dashboard permanecen responsivos; el worker entra en
        \texttt{FAULT} si vence el timeout. & Aprobado \\
        Comandos rápidos &
        El comando pendiente anterior se sustituye por el más reciente. &
        Aprobado \\
        Iluminación, fondo y distancia &
        Se incorporaron tasas sintéticas en E1--E5; requieren confirmación con
        secuencias reales. & Simulado \\
        \bottomrule
    \end{tabularx}
    \caption{Evidencia lógica y simulada de robustez.}
    \label{tab:resultados_robustez}
\end{table}

\begin{table}[H]
    \centering
    \renewcommand{\arraystretch}{1.2}
    \small
    \begin{tabularx}{\textwidth}{p{1.2cm} p{3.3cm} p{2.4cm} p{2.4cm} X}
        \toprule
        \textbf{Esc.} & \textbf{Condiciones} & \textbf{Brazo válido} &
        \textbf{Mano válida} & \textbf{Lectura simulada} \\
        \midrule
        E1 & Alta / media / baja &
        $97.1 / 96.2 / 84.7\%$ & $94.6 / 92.8 / 78.9\%$ &
        La baja iluminación produce la mayor caída visual. \\
        E2 & Fondo simple / complejo &
        $96.8 / 91.4\%$ & $93.5 / 86.2\%$ &
        El fondo complejo afecta más a la asociación de la mano. \\
        E3 & Lenta / normal / rápida &
        $97.0 / 95.9 / 88.6\%$ & $93.8 / 92.1 / 82.7\%$ &
        El movimiento rápido aumenta pérdidas y MAE. \\
        E4 & Visible / oclusión parcial &
        $96.5 / 72.8\%$ & $93.0 / 68.4\%$ &
        Las pérdidas mayores a 750 ms activan el watchdog simulado. \\
        E5 & Corta / media / larga &
        $95.8 / 96.4 / 87.9\%$ & $91.0 / 93.1 / 76.4\%$ &
        La distancia larga degrada especialmente el gripper. \\
        \bottomrule
    \end{tabularx}
    \caption{Tasas sintéticas de detección en los escenarios E1--E5.}
    \label{tab:robustez_simulada}
\end{table}

% ============================================================
\section{Análisis y Discusión de los Resultados}
\label{sec:discusion_resultados}
% ============================================================

Los resultados disponibles permiten afirmar que la arquitectura modular y las
protecciones del simulador funcionan conforme al contrato implementado. En
particular, la validación independiente evita que una mano temporalmente
oculta interrumpa el brazo, y evita que una pose inválida produzca objetivos
extremos. El lote de hombro y codo conserva coherencia a nivel de validación,
aunque las dos RPC se ejecutan secuencialmente y no forman una transacción
atómica dentro de la escena.

La media móvil de cinco muestras reduce fluctuaciones frame a frame, pero
introduce inercia temporal. Con los 27.6 FPS del escenario nominal, su retardo
de grupo aproximado es de 72 ms. Esta estimación proviene de la simulación y
debe medirse durante la campaña real. De forma similar, ejecutar Hands cada dos
frames disminuye la carga de cómputo a cambio de reutilizar temporalmente la
última detección.

El mapeo lineal comprime $180^\circ$ de movimiento humano en $60^\circ$ del
robot. Esta decisión mantiene los objetivos dentro del intervalo configurado,
pero no reproduce la cinemática completa del miembro superior. El UR5 y el
brazo humano poseen estructuras y orientaciones distintas; en esta fase no se
resuelve cinemática inversa ni se intenta conservar la pose cartesiana, a
diferencia de trabajos de transferencia cinemática más completos
~\cite{wang2020fast,han2024robust}.

Frente al Informe \#2, la principal mejora es el paso de una visualización
conceptual a una simulación controlada con estados verificables. El uso de una
cámara monocular mantiene bajo el costo del sistema, en línea con propuestas
de imitación visual con webcam~\cite{xuan2017vision}, pero limita la
profundidad y la precisión de la orientación.

Los resultados sintéticos sugieren que el escenario nominal podría cumplir los
umbrales de 20 FPS, 90\% de detecciones válidas y $5^\circ$ de MAE. La misma
simulación predice incumplimiento bajo iluminación baja, movimiento rápido,
oclusión parcial y distancia larga. Estos patrones son coherentes con las
limitaciones del pipeline, pero dependen de los valores asumidos y no deben
interpretarse como desempeño observado.

Las incidencias encontradas durante el desarrollo fueron relevantes para el
diseño final:

\begin{itemize}
    \item el uso inicial de 11/13/15 seleccionaba el brazo izquierdo cuando la
    inferencia se realizaba sin espejo;
    \item una muestra inválida podía contaminar el suavizado o dejar pendiente
    un objetivo anterior;
    \item el cliente ZeroMQ debía configurar por separado el presupuesto de
    conexión y el timeout operativo;
    \item la escena incluida necesitaba un perfil dinámico compatible para que
    velocidad, aceleración y jerk fueran aplicables; y
    \item consultar aliases después de detectar una colisión podía retrasar o
    impedir la RPC de parada, por lo que se precargaron en el preflight.
\end{itemize}

Dentro de la simulación, las mejores condiciones corresponden a iluminación
alta o media, fondo simple, movimiento lento o normal y distancia intermedia.
Las peores corresponden a iluminación baja, movimiento rápido, oclusión y
distancia larga. El siguiente paso experimental debe ejecutar E1--E5 para
confirmar o refutar esta comparación y reemplazar los valores sintéticos por
FPS, latencia, MAE y porcentajes observados.

% ============================================================
\section{Limitaciones de la Implementación}
\label{sec:limitaciones_implementacion}
% ============================================================

\begin{itemize}
    \item Una cámara RGB monocular no proporciona profundidad métrica; los
    ángulos se estiman en el plano de imagen.
    \item La iluminación, las oclusiones, la ropa y un fondo complejo pueden
    reducir la estabilidad de Pose y Hands.
    \item Los umbrales de apertura del gripper están expresados en píxeles y no
    se normalizan por el tamaño aparente de la mano ni por la distancia a la
    cámara.
    \item Solo se mapean hombro y codo; no se controla la orientación de muñeca
    ni la pose cartesiana completa del efector final.
    \item La detección de colisiones es reactiva después del contacto. No existe
    planificación de trayectorias ni predicción de colisiones.
    \item Las autocolisiones se excluyen deliberadamente; solo se comparan las
    formas móviles del UR5/RG2 con objetos externos al árbol \texttt{/UR5}.
    \item Las RPC de hombro y codo se validan como lote, pero no se ejecutan
    mediante una transacción atómica embebida en la escena.
    \item El retorno a \textit{home} depende de que la comunicación, la
    simulación y el readback permanezcan disponibles durante un máximo de 5 s.
    \item Los FPS, la latencia extremo a extremo, el MAE angular y la robustez
    cuantitativa presentados son simulados; no se midieron con la cámara física.
    \item Las protecciones se diseñaron exclusivamente para CoppeliaSim. El
    software no es un sistema de seguridad certificado y no debe conectarse a
    un robot real sin una arquitectura independiente para hardware.
\end{itemize}
